% ----------------------------------------------------------------------
%  3.3.3. Промежуточные фазы
% ----------------------------------------------------------------------
\subsection{Промежуточные фазы}\label{ssec:intermediate}

\term{Промежуточные фазы}~— все новые фазы двойных и многокомпонентных
систем, которые не прилегают на диаграмме состояния к чистым
компонентам. Их появление часто сопровождается резким изменением
структуры и свойств сплава. Промежуточные фазы, сохраняющие дальний
порядок во всём интервале своей устойчивости, называют
\term{интерметаллическими соединениями}; обычно они существуют при
определённом соотношении компонентов и имеют узкую область
гомогенности.

\paragraph{Что определяет устойчивость промежуточных фаз}
Основные факторы сведены в табл.~\ref{tab:inter-factors}. Как правило,
они действуют совместно: и размерный, и электронный факторы влияют,
например, на фазы Лавеса, $\sigma$-фазы, фазы типа NiAs и CaF$_2$. В
конечном счёте устойчивость фазы определяется термодинамикой. Мы
подробно разберём два «чистых» случая: электронные соединения
(решает электронная концентрация) и фазы внедрения (решает размерный
фактор).

\begin{table}[ht]
\centering
\caption{Факторы, определяющие образование промежуточных фаз}
\label{tab:inter-factors}
\small
\begin{tabularx}{\textwidth}{@{}>{\raggedright\arraybackslash}p{42mm}>{\raggedright\arraybackslash}X@{}}
\toprule
Фактор & Примеры фаз \\
\midrule
различие электроотрицательностей & соединения нормальной валентности (типа NiAs, CaF$_2$) \\
размерный фактор & фазы Лавеса, фазы внедрения \\
электронная концентрация & электронные соединения (фазы Юм-Розери): $\beta$-, $\gamma$-, $\varepsilon$-фазы \\
плотное заполнение пространства (к.\,ч. больше 12, до 16) & фазы Лавеса (MgCu$_2$, MgZn$_2$), $\sigma$-фазы (FeCr, FeV), фазы типа $\beta$-W (A$_3$B), $\delta$ (MoNi), $P$ (MoNiCr), $R$ (MoCrCo), $\mu$ (Mo$_6$Co$_7$), $\chi$ (Re$_{24}$Ti$_5$) \\
\bottomrule
\end{tabularx}
\end{table}

По кристаллической структуре промежуточные фазы делят на структурные
типы, аналогичные чистым металлам (в том числе сложным: $\beta$-U,
$\beta$-W, $\alpha$- и $\beta$-Mn); типы $\gamma$-латуни; типы с высокими
(фазы Лавеса и др.) и низкими (NiAs, CaF$_2$ и др.) координационными
числами; соединения металлов с неметаллами промежуточного характера~—
фазы внедрения.

\paragraph{Электронные соединения}
Мы видели, что растворимость Zn, Ga, Ge в меди ограничена электронной
концентрацией около 1,36; так же ведут себя эти элементы в серебре и
золоте. Оказывается, электронная концентрация определяет и области
существования многих \emph{промежуточных} фаз. Рассмотрим все фазы
системы Cu--Zn (рис.~\ref{fig:CuZn}).
\begin{itemize}
\item $\alpha$~— первичный твёрдый раствор Zn в Cu (ГЦК, 4 атома в
ячейке); предел растворимости задаётся предельной электронной
концентрацией.
\item $\beta$~— фаза около состава CuZn ($\sim$50~ат.\,\% Zn), $e/a=3/2$,
ОЦК (2 атома в ячейке); при низких температурах упорядочивается
($\beta'$).
\item $\gamma$~— около Cu$_5$Zn$_8$, $e/a=21/13$; сложная кубическая
решётка с 52 атомами в ячейке. Её можно представить как 27 малых
ОЦК-ячеек ($3\times3\times3$), из которых удалены два атома~— в вершинах
и в центре большого куба.
\item $\varepsilon$~— около CuZn$_3$, $e/a=7/4$, ГПУ.
\item $\eta$~— первичный твёрдый раствор Cu в Zn (ГПУ, $c/a=1{,}86$).
\end{itemize}
Все промежуточные фазы здесь~— твёрдые растворы с довольно широкими
областями гомогенности и металлической связью.

\begin{figure}[ht]
\centering
\begin{tikzpicture}[x=0.125cm, y=1cm]
  \foreach \a/\b/\col/\lab in {0/38/fmBlue!25/$\alpha$, 38/46/fmGray!12/, 46/50.5/fmAmber!40/$\beta'$,
     50.5/58/fmGray!12/, 58/67/fmGreen!30/$\gamma$, 67/78/fmGray!12/, 78/87/fmRed!25/$\varepsilon$,
     87/97/fmGray!12/, 97/100/fmPurple!25/$\eta$}{
    \fill[\col] (\a,0) rectangle (\b,1);
    \node[font=\small] at ({(\a+\b)/2},0.5) {\lab};}
  \draw[line width=0.5pt] (0,0) rectangle (100,1);
  \foreach \x in {0,20,40,60,80,100}{\draw (\x,0) -- (\x,-0.12) node[below, font=\scriptsize]{\x};}
  \node[font=\small] at (50,-0.75) {Zn, ат.\,\%};
  % электронная концентрация
  \draw[line width=0.5pt] (0,1.6) -- (100,1.6);
  \foreach \c/\t in {0/{1,0}, 36/{1,36}, 50/{3/2}, 61.5/{21/13}, 75/{7/4}, 100/{2,0}}{
    \draw (\c,1.6) -- (\c,1.72) node[above, font=\scriptsize]{\t};}
  \node[font=\small, anchor=east] at (-2,1.6) {$e/a$};
  \node[font=\scriptsize, fmGray] at (42,0.5) {$\alpha{+}\beta'$};
  \node[font=\scriptsize, fmGray] at (54.2,0.5) {$\beta'{+}\gamma$};
  \node[font=\scriptsize, fmGray] at (72.5,0.5) {$\gamma{+}\varepsilon$};
\end{tikzpicture}
\caption{Фазы системы Cu--Zn вдоль оси составов при $\sim$300--400\,°C
(упрощённо) и соответствующая электронная концентрация $e/a=1+c$.
Границы областей $\alpha$, $\beta'$, $\gamma$, $\varepsilon$, $\eta$
расположены около характерных значений $e/a$}
\label{fig:CuZn}
\end{figure}

\begin{opredelenie}
\term{Электронные соединения} (фазы Юм-Розери)~— металлические фазы,
структура которых определяется концентрацией валентных электронов, а не
обычными правилами валентности. Их образуют металлы двух классов:
1)~металлы группы меди и переходные металлы; 2)~$B$-металлы II--V групп.
\end{opredelenie}

При подсчёте электронной концентрации число валентных электронов
принимают равным номеру группы; переходным металлам (из-за
незаполненной $d$-оболочки и перекрытия уровней $(n-1)d$ и $ns$)
приписывают чаще всего нулевую валентность, реже 1 или~2. Примеры
(табл.~\ref{tab:electron-comp}): у Cu$_3$Al $e/a=(3\cdot1+3)/4=3/2$, у
Cu$_5$Sn $(5+4)/6=3/2$, у NiAl $(0+3)/2=3/2$.

\begin{table}[ht]
\centering
\caption{Электронные соединения}
\label{tab:electron-comp}
\small
\begin{tabularx}{\textwidth}{@{}llc>{\raggedright\arraybackslash}X@{}}
\toprule
Фаза & Структура & $e/a$ & Примеры \\
\midrule
$\beta$ & ОЦК (иногда ГПУ, $\beta$-Mn) & $3/2=1{,}50$ & CuZn, AgMg, AgZn, AuCd, AuMg, AuZn, Cu$_3$Al, Cu$_5$Sn, NiAl \\
$\gamma$ & кубическая, 52 атома & $21/13=1{,}62$ & Cu$_5$Zn$_8$, Ag$_5$Cd$_8$, Ag$_5$Hg$_8$, Cu$_{31}$Sn$_8$, Na$_{31}$Pb$_8$ \\
$\varepsilon$ & ГПУ & $7/4=1{,}75$ & CuZn$_3$, AgCd$_3$, Au$_5$Al$_3$, Cu$_3$Si, Cu$_3$Sn \\
\bottomrule
\end{tabularx}
\end{table}
\noindent В оригинале среди примеров указаны Ni$_{31}$Pb$_8$ и Ni$_3$Sn,
которые правилу не подчиняются; по-видимому, имелись в виду
Na$_{31}$Pb$_8$ и Cu$_3$Sn\ispr{Для Ni$_{31}$Pb$_8$ и Ni$_3$Sn при нулевой
валентности никеля получается $e/a=32/39$ и $4/4$, а не $21/13$ и $7/4$.
У Na$_{31}$Pb$_8$: $(31+32)/39=21/13$; у Cu$_3$Sn: $(3+4)/4=7/4$.}.

\paragraph{Почему важна электронная концентрация: зонное объяснение}
Объяснение дал Джонс на основе зонной теории (п.~\ref{ssec:bands}).
\begin{shagi}
\item С ростом электронной концентрации сфера Ферми растёт. Пока она
далеко от границ зоны Бриллюэна, плотность состояний $g(E)$ следует
параболе.
\item Когда сфера приближается к границе, $g(E)$ растёт быстрее параболы
и в момент касания достигает пика (участок $AB$ на
рис.~\ref{fig:dos-band}). Затем $g(E)$ резко падает.
\item Пока $g(E)$ велика, новые электроны размещаются «дёшево»~— на
близко лежащих уровнях. После пика каждый следующий электрон
обходится дороже, и энергия решётки растёт быстрее.
\item Для ГЦК касание происходит при $e/a=1{,}36$, для ОЦК~— при
$1{,}48$ (п.~\ref{ssec:bands}). Поэтому при $e/a>1{,}36$ ГЦК-решётка
теряет выгодность, и при $e/a$ около 1,48 устойчивее ОЦК-фаза
(рис.~\ref{fig:jones}). В промежутке существует смесь $\alpha$- и
$\beta$-фаз.
\end{shagi}
Аналогично объясняются переходы $\beta\to\gamma$ и
$\gamma\to\varepsilon$: касание поверхности Ферми с границей зоны для
$\gamma$-латуни происходит при $e/a=1{,}54$, для $\varepsilon$-фазы~— при
$1{,}75$.

\begin{figure}[ht]
\centering
\begin{tikzpicture}
\begin{axis}[width=0.66\textwidth, height=5.4cm, xmin=0, xmax=1.0, ymin=0, ymax=1.5,
  xlabel={$E$}, ylabel={$g(E)$}, xtick={0.55,0.7}, xticklabels={$e/a=1{,}36$,$1{,}48$},
  ytick=\empty, axis lines=left, xticklabel style={font=\scriptsize}]
  \addplot[fmGray, dashed, domain=0:1, samples=50] {1.05*sqrt(x)};
  \addplot[fmBlue, thick, smooth] coordinates {(0,0) (0.1,0.33) (0.3,0.58) (0.45,0.76) (0.52,0.92) (0.55,1.12) (0.57,0.82) (0.65,0.55) (0.8,0.42) (1.0,0.35)};
  \addplot[fmRed, thick, smooth] coordinates {(0,0) (0.1,0.33) (0.3,0.58) (0.5,0.75) (0.62,0.9) (0.68,1.08) (0.70,1.3) (0.72,0.95) (0.8,0.62) (1.0,0.45)};
  \node[font=\small, fmBlue] at (axis cs:0.45,1.15) {ГЦК};
  \node[font=\small, fmRed] at (axis cs:0.82,1.2) {ОЦК};
\end{axis}
\end{tikzpicture}
\caption{Плотность состояний ГЦК- и ОЦК-фаз (схема по Джонсу). Пики
отвечают касанию сферы Ферми с границей зоны Бриллюэна: для ГЦК при
$e/a=1{,}36$, для ОЦК при $1{,}48$. После пика ГЦК-кривая падает, и
дальнейшее заполнение выгоднее в ОЦК-решётке~— так объясняется граница
$\alpha$- и $\beta$-фаз латуни}
\label{fig:jones}
\end{figure}

\paragraph{Фазы внедрения}
\term{Фазы внедрения}~— соединения переходных металлов с водородом,
углеродом, азотом (реже с кислородом), строение которых определяется
\emph{размерным} фактором. По правилу Хэгга:
\begin{itemize}
\item если $r_X/r_\text{Me}\le0{,}59$, атомы металла образуют
простую решётку (ГЦК, ОЦК, ГПУ, иногда простую гексагональную с
к.\,ч.~8 и $c/a\approx1$), а атомы неметалла $X$ сидят в её пустотах;
\item если $r_X/r_\text{Me}>0{,}59$, образуются соединения со сложными
решётками, уже не являющиеся фазами внедрения.
\end{itemize}

Почти все фазы внедрения имеют решётку, \emph{отличную} от решётки
исходного металла. Число внедрённых атомов сравнимо с числом атомов
металла, решётка сильно искажается, возникает как бы «внутреннее
давление», и металлическая подрешётка перестраивается~— словно
претерпевает полиморфное превращение.

Я.\,С.~Уманский предположил, что H, C и N в фазах внедрения отдают часть
своих электронов в перекрывающиеся $(n-1)d$- и $ns$-зоны переходного
металла и находятся как бы в металлическом состоянии. В пользу этого
говорят металлические свойства фаз внедрения: металлический блеск,
хорошая электропроводность (у ZrN и TiN даже выше, чем у чистых
металлов), снижение сопротивления при охлаждении. Модель объясняет и
широкие области гомогенности, несвойственные обычным ионным (NaCl) или
ковалентным (SiC) соединениям. Фазы внедрения можно рассматривать как
твёрдые растворы C, N, H в металлах, испытавших полиморфное
превращение; переменный состав (например, у карбида титана TiC) тогда
не требует особых объяснений.

Типичные составы~— MeX, Me$_2$X, Me$_4$X, MeX$_2$ (рис.~\ref{fig:interstitial}):
\begin{itemize}
\item \emph{MeX}: в карбидах и нитридах Ti, Zr, Hf, V, Nb, Ta атомы
металла образуют ГЦК-решётку, а C или N занимают октаэдрические пустоты~—
структура NaCl, но не все позиции заполнены; в WC атомы металла образуют
простую гексагональную решётку;
\item \emph{Me$_4$X} (нитрид железа Fe$_4$N): ГЦК-решётка железа, азот
занимает четверть октаэдрических пустот~— центр ячейки;
\item \emph{Me$_2$X} (W$_2$C): атомы углерода~— примерно в половине
октаэдрических пустот ГПУ-решётки вольфрама.
\end{itemize}

\begin{figure}[ht]
\centering
\begin{subfigure}[b]{0.32\textwidth}\centering
\begin{tikzpicture}[cubeproj, scale=1.85, anion/.style={atom=4.6mm}, cation/.style={atomI=2.4mm}]
  \CubeEdges[draw opacity=0]
  % сгенерировано tools/gen_ionic_cells.py — не править вручную
\draw[bondline] (0,0,0) -- (0,0,0.5);
\draw[bondline] (0,0,0) -- (0,0.5,0);
\draw[bondline] (0,0,0) -- (0.5,0,0);
\draw[bondline] (0,0,0.5) -- (0,0,1);
\draw[bondline] (0,0,0.5) -- (0,0.5,0.5);
\draw[bondline] (0,0,0.5) -- (0.5,0,0.5);
\draw[bondline] (0,0,1) -- (0,0.5,1);
\draw[bondline] (0,0,1) -- (0.5,0,1);
\draw[bondline] (0,0.5,0) -- (0,0.5,0.5);
\draw[bondline] (0,0.5,0) -- (0,1,0);
\draw[bondline] (0,0.5,0) -- (0.5,0.5,0);
\draw[bondline] (0,0.5,0.5) -- (0,0.5,1);
\draw[bondline] (0,0.5,0.5) -- (0,1,0.5);
\draw[bondline] (0,0.5,0.5) -- (0.5,0.5,0.5);
\draw[bondline] (0,0.5,1) -- (0,1,1);
\draw[bondline] (0,0.5,1) -- (0.5,0.5,1);
\draw[bondline] (0,1,0) -- (0,1,0.5);
\draw[bondline] (0,1,0) -- (0.5,1,0);
\draw[bondline] (0,1,0.5) -- (0,1,1);
\draw[bondline] (0,1,0.5) -- (0.5,1,0.5);
\draw[bondline] (0,1,1) -- (0.5,1,1);
\draw[bondline] (0.5,0,0) -- (0.5,0,0.5);
\draw[bondline] (0.5,0,0) -- (0.5,0.5,0);
\draw[bondline] (0.5,0,0) -- (1,0,0);
\draw[bondline] (0.5,0,0.5) -- (0.5,0,1);
\draw[bondline] (0.5,0,0.5) -- (0.5,0.5,0.5);
\draw[bondline] (0.5,0,0.5) -- (1,0,0.5);
\draw[bondline] (0.5,0,1) -- (0.5,0.5,1);
\draw[bondline] (0.5,0,1) -- (1,0,1);
\draw[bondline] (0.5,0.5,0) -- (0.5,0.5,0.5);
\draw[bondline] (0.5,0.5,0) -- (0.5,1,0);
\draw[bondline] (0.5,0.5,0) -- (1,0.5,0);
\draw[bondline] (0.5,0.5,0.5) -- (0.5,0.5,1);
\draw[bondline] (0.5,0.5,0.5) -- (0.5,1,0.5);
\draw[bondline] (0.5,0.5,0.5) -- (1,0.5,0.5);
\draw[bondline] (0.5,0.5,1) -- (0.5,1,1);
\draw[bondline] (0.5,0.5,1) -- (1,0.5,1);
\draw[bondline] (0.5,1,0) -- (0.5,1,0.5);
\draw[bondline] (0.5,1,0) -- (1,1,0);
\draw[bondline] (0.5,1,0.5) -- (0.5,1,1);
\draw[bondline] (0.5,1,0.5) -- (1,1,0.5);
\draw[bondline] (0.5,1,1) -- (1,1,1);
\draw[bondline] (1,0,0) -- (1,0,0.5);
\draw[bondline] (1,0,0) -- (1,0.5,0);
\draw[bondline] (1,0,0.5) -- (1,0,1);
\draw[bondline] (1,0,0.5) -- (1,0.5,0.5);
\draw[bondline] (1,0,1) -- (1,0.5,1);
\draw[bondline] (1,0.5,0) -- (1,0.5,0.5);
\draw[bondline] (1,0.5,0) -- (1,1,0);
\draw[bondline] (1,0.5,0.5) -- (1,0.5,1);
\draw[bondline] (1,0.5,0.5) -- (1,1,0.5);
\draw[bondline] (1,0.5,1) -- (1,1,1);
\draw[bondline] (1,1,0) -- (1,1,0.5);
\draw[bondline] (1,1,0.5) -- (1,1,1);
\node[anion] at (0,1,0) {};
\node[cation] at (0.5,1,0) {};
\node[anion] at (1,1,0) {};
\node[cation] at (0,1,0.5) {};
\node[anion] at (0.5,1,0.5) {};
\node[cation] at (1,1,0.5) {};
\node[anion] at (0,1,1) {};
\node[cation] at (0.5,1,1) {};
\node[anion] at (1,1,1) {};
\node[cation] at (0,0.5,0) {};
\node[anion] at (0.5,0.5,0) {};
\node[cation] at (1,0.5,0) {};
\node[anion] at (0,0.5,0.5) {};
\node[cation] at (0.5,0.5,0.5) {};
\node[anion] at (1,0.5,0.5) {};
\node[cation] at (0,0.5,1) {};
\node[anion] at (0.5,0.5,1) {};
\node[cation] at (1,0.5,1) {};
\node[anion] at (0,0,0) {};
\node[cation] at (0.5,0,0) {};
\node[anion] at (1,0,0) {};
\node[cation] at (0,0,0.5) {};
\node[anion] at (0.5,0,0.5) {};
\node[cation] at (1,0,0.5) {};
\node[anion] at (0,0,1) {};
\node[cation] at (0.5,0,1) {};
\node[anion] at (1,0,1) {};

\end{tikzpicture}
\caption{TiC, ZrC, NbN (MeX)}
\end{subfigure}\hfill
\begin{subfigure}[b]{0.32\textwidth}\centering
\begin{tikzpicture}[cubeproj, scale=1.85]
  \CubeEdges
  \FCCback[4.4mm]
  \node[atomI=2.6mm] at (0.5,0.5,0.5) {};
  \FCCmid[4.4mm]
  \FCCfront[4.4mm]
\end{tikzpicture}
\caption{Fe$_4$N (Me$_4$X)}
\end{subfigure}\hfill
\begin{subfigure}[b]{0.32\textwidth}\centering
\begin{tikzpicture}[cubeproj, scale=1.85]
  \CubeEdges
  \FCCback[4.4mm]
  \foreach \p in {(0.5,1,0),(0.5,1,1),(0,1,0.5),(1,1,0.5)}{\node[vac=2.2mm] at \p {};}
  \foreach \p in {(0,0.5,0),(1,0.5,0),(0,0.5,1),(1,0.5,1)}{\node[vac=2.2mm] at \p {};}
  \node[atomI=2.6mm] at (0.5,0.5,0.5) {};
  \FCCmid[4.4mm]
  \foreach \p in {(0.5,0,0),(0.5,0,1),(0,0,0.5),(1,0,0.5)}{\node[vac=2.2mm] at \p {};}
  \FCCfront[4.4mm]
\end{tikzpicture}
\caption{аустенит: октапустоты}
\end{subfigure}
\caption{Атомы неметалла (красные) в октаэдрических пустотах ГЦК-решётки
металла: \textit{а}~— все пустоты заняты (структура NaCl);
\textit{б}~— занята одна из четырёх (центр ячейки); \textit{в}~—
все октаэдрические пустоты ГЦК-ячейки (центр куба и середины рёбер,
пунктир), которые может занимать углерод в аустените}
\label{fig:interstitial}
\end{figure}

\paragraph{Какие пустоты занимает неметалл}
Действует правило соотношения размеров (п.~\ref{sssec:ionic}). При
$r_X/r_\text{Me}=0{,}41$--$0{,}59$ неметалл сидит в октаэдрических
пустотах. При отношении меньше 0,41 одиночный атом «болтается» в
октаэдрической пустоте, и решётка неустойчива. Выходов два: в одну
октаэдрическую пустоту садятся \emph{два} атома неметалла, или одиночные
атомы переходят в тетраэдрические пустоты. Пример второго случая~—
гидрид циркония: $r_\text{H}/r_\text{Zr}=0{,}046/0{,}160=0{,}29$, и
водород занимает тетраэдрические пустоты.

\paragraph{Свойства фаз внедрения}
\begin{itemize}
\item Фазы внедрения с однотипными решётками, как правило, образуют между
собой непрерывные ряды твёрдых растворов замещения (NbC--TaC, NbC--TiC,
NbC--ZrC, TaC--TiC, TaC--ZrC, TiC--VC; нитриды ZrN, NbN, TiN, VN, TaN).
Исключения~— ZrC--VC и ZrN--VN: радиус атома Zr на 21\,\% больше, чем
V, и размерный фактор неблагоприятен.
\item Очень высокие температуры плавления, часто выше, чем у исходных
металлов, и твёрдость, приближающаяся к твёрдости алмаза. Поэтому
считают, что связь металл--неметалл в них частично ковалентная.
\item Многие фазы внедрения при низких температурах становятся
сверхпроводниками (NbN~— около 16~К)\ispr{В оригинале (с.~420)
сказано, что сверхпроводимостью обладают «почти все» фазы внедрения;
это преувеличение: сверхпроводящими являются многие (NbN, NbC, MoC, TaC
и др.), но не все.}.
\end{itemize}

\begin{itogi}
\begin{itemize}
\item Твёрдые растворы бывают замещения, внедрения и вычитания; тип
устанавливают, сравнивая число атомов в ячейке $n=\rho V/(m_uA_\text{ср})$
с числом $n_0$ для чистой основы.
\item Растворение искажает решётку: растут сопротивление и твёрдость,
период меняется (закон Вегарда~— лишь приближение).
\item Неограниченная растворимость требует одинаковых решёток, разницы
диаметров не более 14--15\,\% и близких электрохимических свойств; эти
условия необходимы, но не достаточны.
\item Растворимость $B$-металлов в меди ограничена электронной
концентрацией около 1,36~— моментом касания сферы Ферми с границей зоны
Бриллюэна ГЦК-решётки.
\item Упорядочение (сверхструктуры Cu$_3$Au, CuAu, CuZn) снижает
искажения и электросопротивление и даёт сверхструктурные линии.
\item Электронные соединения: $\beta$ ($e/a=3/2$), $\gamma$ ($21/13$),
$\varepsilon$ ($7/4$). Фазы внедрения ($r_X/r_\text{Me}\le0{,}59$)~—
металлоподобные, очень твёрдые и тугоплавкие.
\end{itemize}
\end{itogi}

\begin{proverka}
\item Чем раствор вычитания отличается от раствора внедрения? Как их
различить экспериментально?
\item Почему литий не смешивается с калием даже в жидком состоянии?
\item Рассчитайте предельную растворимость алюминия в меди, если она
ограничена электронной концентрацией 1,4.
\item Почему у сверхструктуры Ni$_3$Fe не видны сверхструктурные линии на
рентгенограмме?
\item Чему равна электронная концентрация в Cu$_9$Al$_4$? К какому типу
электронных соединений его отнести?
\item Почему углерод в $\gamma$-Fe занимает октаэдрические, а не
тетраэдрические пустоты?
\end{proverka}
